% Final technical spec sheet (BGU \S 10.7.2.5).
\section{Final Technical Specification}
\label{sec:spec}

This section reports the as-built specification of the delivered software
product, and documents every drift from the preliminary specification
\cite{shi2024transnext, huang2017densenet}.

\subsection{Product overview}
The delivered product is a research-grade Python package that:
\begin{enumerate}
  \item synthetically degrades a clean image dataset using a deterministic
        five-axis pipeline (saturation, downsample, blur, additive noise,
        salt and pepper);
  \item fine-tunes one of three pretrained deep backbones (ResNet50,
        DenseNet121, TransNeXt-tiny) on the degraded distribution under a
        unified training recipe;
  \item evaluates the trained model on a held-out validation split and
        records validation accuracy, calibration, inference throughput, and
        per-pipeline PSNR / SSIM metrics;
  \item aggregates results across the 402-cell campaign and the 48-replicate
        multi-seed audit into a single canonical JSON artifact and a 7-tab
        interactive dashboard.
\end{enumerate}

\subsection{Inputs, outputs, interfaces}
\textbf{Inputs.}\quad Standard image datasets (CIFAR-10 \cite{krizhevsky2009cifar10}
and MNIST \cite{lecun1998mnist}), automatically downloaded by torchvision on
first invocation. A YAML / dictionary configuration object selects the
backbone, the degradation phase, the level, the axis (for single-axis
phases), the treatment (for Phase D), and the random seed.

\textbf{Outputs.}\quad Per-cell training logs, validation metrics, best
checkpoint, image-quality probe JSON. The canonical aggregator
(\texttt{Final\_Exp.json}, schema version 3) joins the 402 canonical cells
with the 48 multi-seed replicates and exposes
$\langle\mathrm{val\_acc}\rangle \pm \sigma$ pairs.

\textbf{Hardware interface.}\quad The training stack runs on a single NVIDIA
RTX 5070 (sm\_120, Blackwell architecture) at \texttt{bf16-mixed} precision.
A CPU-only fallback (\texttt{32-true} precision) is available for portability
but is too slow to be useful at the campaign scale. There is no specialized
hardware development; the system targets commodity GPU servers.

\subsection{Final specification of the system}
\textbf{Models.}\quad Three backbones are supported. Table \ref{tab:hparams}
summarizes the training recipe; the differential learning-rate convention
(head $\times 10$ over backbone) is taken from the source papers.

\begin{table}[htbp]
\centering
\caption{Training hyperparameters (final).}
\label{tab:hparams}
\setlength{\tabcolsep}{4pt}
\renewcommand{\arraystretch}{1.1}
\begin{tabular}{lll}
\toprule
Parameter & CNNs (ResNet50, DenseNet121) & TransNeXt-tiny \\
\midrule
Input resolution        & $224 \times 224$        & $224 \times 224$ \\
Optimizer               & AdamW \cite{loshchilov2019adamw} & AdamW \\
Schedule                & Cosine annealing        & Cosine annealing \\
Head learning rate      & $10^{-3}$               & $5 \times 10^{-4}$ \\
Backbone learning rate  & $10^{-4}$               & $5 \times 10^{-5}$ \\
Weight decay            & $10^{-4}$               & $5 \times 10^{-2}$ \\
Label smoothing         & 0.1                     & 0.1 \\
Stochastic depth        & 0.0                     & 0.1 \\
Gradient clip           & $\|\nabla\|_2 \leq 1.0$ & $\|\nabla\|_2 \leq 1.0$ \\
Batch size              & 32                      & 32 \\
Max epochs              & 60                      & 60 \\
Early stopping          & patience 10             & patience 10 \\
Precision               & \texttt{bf16-mixed}     & \texttt{bf16-mixed} \\
Seed                    & 42 (canonical)          & 42 (canonical) \\
\bottomrule
\end{tabular}
\end{table}

\textbf{Degradation pipeline.}\quad The five-level schedule (Table
\ref{tab:levels}) is the single source of truth for L1 through L5; it is
imported from \texttt{src/data/degradation\_levels.py} into both the training
and the evaluation paths.

\begin{table}[htbp]
\centering
\caption{Five-level synthetic degradation schedule. Blur kernel and $\sigma$ are pixel-domain values at the $224 \times 224$ model input.}
\label{tab:levels}
\setlength{\tabcolsep}{4pt}
\renewcommand{\arraystretch}{1.1}
\begin{tabular}{llrrrrrr}
\toprule
Level & Name & low\_res & Blur K & Blur $\sigma$ & Noise std & S\&P & Saturation \\
\midrule
L1 & Mild & 18 & 13 & 2.50 & 0.04 & 0.03 & 0.95 \\
L2 & Light & 12 & 25 & 5.00 & 0.08 & 0.06 & 0.65 \\
L3 & Moderate & 8 & 41 & 8.00 & 0.12 & 0.10 & 0.40 \\
L4 & Severe & 6 & 61 & 12.00 & 0.16 & 0.14 & 0.15 \\
L5 & Extreme & 3 & 91 & 18.00 & 0.22 & 0.18 & 0.00 \\
\bottomrule
\end{tabular}
\end{table}


\textbf{Data pipeline.}\quad Pipeline version 2 (US-017, finalized
2026-05-20): the saturation lerp is applied first, then the bilinear
downsample to \texttt{low\_res}, then additive Gaussian noise and salt and
pepper at the low resolution, then bicubic upsample to $224 \times 224$,
then Gaussian blur (kernel calibrated for $224 \times 224$), then a final
clip to $[0, 1]$ and ImageNet normalization.

\textbf{Determinism.}\quad The per-sample degradation seed is
$\text{idx} + 10^7$ on the validation split and \texttt{idx} on the training
split. The byte-identity gate in
\texttt{src/tests/test\_degradation\_determinism.py} asserts that two
independent reads of the same validation index yield bit-identical pixels
across freshly built datasets; this is also verified across seeds 42, 43,
and 44 by the multi-seed audit test.

\subsection{Drift from the preliminary specification}

The Preliminary Design Review (PDR, \texttt{Pre-2026-061.pdf}) defined the
project scope. Five drifts occurred during implementation:

\begin{enumerate}
  \item \textbf{TransNeXt variant: small to tiny.} The preliminary spec
        proposed \texttt{transnext\_small} ($\sim 50$M parameters); US-016
        (2026-05-14) substituted \texttt{transnext\_tiny} ($\sim 28$M
        parameters) to keep the parameter count comparable to ResNet50
        ($\sim 25$M) and to recover $\sim 310$ GPU-hours across the
        TransNeXt rows. The change preserves the foveal-attention
        mechanism, which was the architectural property of interest, and it
        matches the published TransNeXt $224 \times 224$ checkpoint exactly.
  \item \textbf{Uniform $224 \times 224$ protocol.} A short-lived V3
        refactor (2026-05-08 through 2026-05-13) attempted native-resolution
        training ($32 \times 32$ for TransNeXt on CIFAR-10, $28 \rightarrow 32$
        for MNIST). The refactor violated the fair-comparison invariant
        (different backbones saw different tensor shapes), so it was rolled
        back. Every cell now trains at $224 \times 224$.
  \item \textbf{Blur kernel rescaling.} The 2026-05-12 degradation curve
        used blur kernels designed for native $32 \times 32$ imagery
        (K $= 3 \ldots 13$, $\sigma = 0.8 \ldots 2.3$). At $224 \times 224$
        those kernels covered $\sim 1$ to $6 \%$ of the image width and
        produced essentially invisible blur. They were rescaled on
        2026-05-14 to pixel-domain values at $224 \times 224$ (K $= 13 \ldots 91$,
        $\sigma = 2.5 \ldots 18$), with the kernel-to-$\sigma$ rule
        $K = 2 \lceil 2.5 \sigma \rceil + 1$ so the kernel always captures the
        $2.5 \sigma$ Gaussian footprint.
  \item \textbf{Optuna freeze.} The preliminary plan reserved a per-(phase,
        cell) hyperparameter sweep. The final spec freezes the Optuna L3
        winner per (model, dataset) and reuses it across every level and
        every axis. The change saves $\sim 200$ GPU-hours and removes a
        per-phase confound from the cross-phase comparison.
  \item \textbf{Multi-seed audit.} The preliminary spec used a single seed
        (42). US-042 (2026-06-01) added a 48-replicate audit at seeds 43
        and 44 across the 24 L3 headline cells so that the published L3
        accuracies carry an explicit $\sigma$ bound rather than a single
        point estimate.
\end{enumerate}

\subsection{Phases of the experimental program}

The 402-cell campaign is partitioned into seven phases. Table
\ref{tab:phase_overview} summarizes each phase.

\begin{table}[htbp]
\centering
\caption{Experimental phases of the 402-cell campaign. Multi-seed audit replicates are layered on top of the 24 L3 headline cells.}
\label{tab:phase_overview}
\setlength{\tabcolsep}{4pt}
\renewcommand{\arraystretch}{1.1}
\begin{tabular}{llrl}
\toprule
Phase & Purpose & Cells & Tag prefix \\
\midrule
A      & Clean baseline (no degradation)                    & 6   & \texttt{final\_clean\_} \\
B      & Combined degradation, all axes at level L         & 30  & \texttt{final\_B\_L*} \\
C      & Single-axis isolation, others at identity         & 150 & \texttt{final\_C\_L*\_<axis>\_} \\
D      & Regularization recovery (T1 / T2 / T3) on Phase B & 90  & \texttt{final\_D\_T*\_L*\_} \\
B2     & THz protocol simplification (sat $=0$, noise $=0$, T3) & 30 & \texttt{final\_B2\_L*\_} \\
B2nr   & B2 without regularization, L3 only                & 6   & \texttt{final\_B2nr\_L3\_} \\
C2     & Single-axis under THz protocol                    & 90  & \texttt{final\_C2\_L*\_<axis>\_} \\
\midrule
Audit  & Multi-seed audit (seeds $\{43, 44\}$) on the 24 L3 headline cells & 48 & \texttt{*\_seed\{43,44\}} \\
\midrule
\textbf{Total} & & \textbf{450} & \\
\bottomrule
\end{tabular}
\end{table}
